\documentclass[12pt,a4paper]{scrartcl}
\usepackage[utf8]{inputenc}
\usepackage[T1]{fontenc}
\usepackage[ngerman]{babel}
\usepackage[a4paper,scale=0.8]{geometry} 
\usepackage{fancyhdr}

\pagestyle{fancy} %eigener Seitenstil
\fancyhf{} %alle Kopf- und Fußzeilenfelder bereinigen
\fancyhead[L]{Philipp Beck} %Kopfzeile links
\fancyhead[C]{Shiftinvariante Unterräume} %zentrierte Kopfzeile
\fancyhead[R]{29.06.2018} %Kopfzeile rechts
\renewcommand{\headrulewidth}{0.4pt} %obere Trennlinie
\fancyfoot[C]{\thepage} %Seitennummer
\renewcommand{\footrulewidth}{0.4pt} %untere Trennlinie
\usepackage{paralist}
\usepackage{enumitem} 

%\usepackage {picins}
\usepackage{color}
\usepackage{float}

% -- Mathe
%\usepackage{gensymb}
%fancyhdr, lastpage, booktabs, xy
\usepackage{graphicx}
\usepackage{amsmath}
\usepackage{amsfonts}
\usepackage{amssymb}

\usepackage{amsthm}

\usepackage{hyperref}
\usepackage{mathrsfs}

% -- Malen
\usepackage{tikz}

\usetikzlibrary{patterns, decorations.pathreplacing, decorations.pathmorphing, arrows}

\usepackage{units}

\numberwithin{equation}{section}

\author{Philipp Beck}

\subtitle{Sommersemester 2018}
\title{Seminar zur mikrolokalen Analysis}

%\makeatletter
%\renewcommand*{\thesection}{\arabic{section}}
%\renewcommand*{\p@section}{\thechapter.}

%Befehle für bestimmte mathematische Symbole
\newcommand{\Gl}{\operatorname{GL}}
\newcommand{\Sl}{\operatorname{SL}}
\newcommand{\Sym}{\operatorname{Sym}}
\newcommand{\ord}{\operatorname{o}}
\newcommand{\id}{\operatorname{id}}
\newcommand{\Aut}{\operatorname{Aut}}
\newcommand{\nt}{\unlhd}
\newcommand{\td}[1]{ \ \text{\rmfamily d} #1}
\newcommand{\Stab}{\operatorname{Stab}}
\newcommand{\Ker}{\operatorname{Ker}}
\newcommand{\Bild}{\operatorname{Im}}
\newcommand{\sgn}{\operatorname{sgn}}
\newcommand{\ggT}{\operatorname{ggT}}
\newcommand{\Syl}{\operatorname{Syl}}
\newcommand{\Grad}{\operatorname{Grad}}
\newcommand{\Char}{\operatorname{char}}
\newcommand{\Quot}{\operatorname{Quot}}
\newcommand{\Hom}{\operatorname{Hom}}
\newcommand{\Gal}{\operatorname{Gal}}
\newcommand{\kgv}{\operatorname{kgV}}
%Mengensymbole
\newcommand{\N}{\mathbb{N}}
\newcommand{\Z}{\mathbb{Z}}
\newcommand{\Zn}[1]{\mathbb{Z}/#1\mathbb{Z}}
\newcommand{\Q}{\mathbb{Q}}
\newcommand{\R}{\mathbb{R}}
\newcommand{\C}{\mathbb{C}}
\renewcommand{\L}{\mathrm{L}}
\newcommand{\F}{\mathbb{F}}
\renewcommand{\i}{\textrm{i}}
\renewcommand{\H}{\mathrm{H}}
\newcommand{\lin}{\operatorname{lin}}
\renewcommand\qedsymbol{$\blacksquare$}

%Big and italic
\newcommand{\bi}[1]{\textbf{\textit{#1}}}

\swapnumbers

\newtheoremstyle{satz}% name
{15pt}% Space above
{5pt}% Space below
{\itshape}% Body font
{}% Indent amount: Indent amount: empty = no indent, \parindent = normal paragraph indent
{\bfseries}% Theorem head font
{:}% Punctuation after theorem head
{\newline}% Space after theorem head: Space after theorem head: { } = normal interword space; \newline = linebreak
{}% Theorem head spec (can be left empty, meaning `normal')

\newtheoremstyle{def}% name
{15pt}% Space above
{5pt}% Space below
{}% Body font
{}% Indent amount: Indent amount: empty = no indent, \parindent = normal paragraph indent
{\bfseries}% Theorem head font
{:}% Punctuation after theorem head
{\newline}% Space after theorem head: Space after theorem head: { } = normal interword space; \newline = linebreak
{}% Theorem head spec (can be left empty, meaning `normal')

\newtheoremstyle{exercise}% name
{15pt}% Space above
{5pt}% Space below
{}% Body font
{}% Indent amount: Indent amount: empty = no indent, \parindent = normal paragraph indent
{\bfseries}% Theorem head font
{:}% Punctuation after theorem head
{0.5em}% Space after theorem head: Space after theorem head: { } = normal interword space; \newline = linebreak
{}% Theorem head spec (can be left empty, meaning `normal')





\newcommand{\thistheoremname}{}
\theoremstyle{satz}
%Satz
\newtheorem{sz}{Satz}[section]

%Lemma
\newtheorem{lemma}[sz]{Lemma}
\newtheorem{gsz}[sz]{\thistheoremname}
\newtheorem*{gthm_no_num}{\thistheoremname}


\theoremstyle{def}
\newtheorem{df}[sz]{Definition}
\newtheorem{gd}[sz]{\thistheoremname}

\newtheorem*{gdf_no_num}{\thistheoremname}


\theoremstyle{exercise}
\newtheorem{exe}{Aufgabe}[section]
\newtheorem{loes}{Lösung}[section]

%Für Beispiele , Bemerkungen
%Also alles was wie eine Definition aussehen soll, aber keine ist.
%\begin{genericdf}{(Hier den Text einfügen)}
%	Inhalt...
%\end{genericdf}
\newenvironment{genericdf}[1]{\renewcommand{\thistheoremname}{#1}\begin{gd}}{\end{gd}}
%Dasselbe ohne Nummerierung
\newenvironment{generic_no_num}[1]{\renewcommand{\thistheoremname}{#1}\begin{gdf_no_num}}{\end{gdf_no_num}}

%Hier für Sätze
\newenvironment{genericthm}[1]{\renewcommand{\thistheoremname}{#1}\begin{gsz}}{\end{gsz}}
%Analog zu Definitionen
\newenvironment{genericthm_no_num}[1]{\renewcommand{\thistheoremname}{#1}\begin{gthm_no_num}}{\end{gthm_no_num}}
\parindent0pt
\parskip1ex


%----Dokument----

\begin{document}
\section{Beurling-Helson und Folgerungen}

\begin{genericthm}{Beurling-Helson}\label{skript:1.1}
Sei $E$ ein $\mathscr{S}$-invarianter Unterraum von $\L^2$.
Dann gibt es folgende zwei Fälle:
\begin{enumerate}[label=\textbf{(\arabic*)}]
\item
Es gilt $\mathscr{S} E = E$.
Dann existiert eine $\lambda_N$-messbares $e \subset \mathbb{T}$ mit $E = \chi_e \L^2$.
\item
Es gilt $\mathscr{S} E \neq E$.
Dann existiert ein $\lambda_N$-messbares $\Theta : \mathbb{T} \to \C$ mit $|\Theta| = 1$ fast überall, sodass
$E = \Theta \H^2$ gilt.
\end{enumerate}
\end{genericthm}

\begin{proof}
\begin{enumerate}[label=\textbf{(\arabic*)}]
\item
Sei $P_E$ die orthogonale Projektion auf $E$ und $\varphi := P_E 1 $.
Das ist $1- \varphi$ wegen $\varphi E = E$ orthogonal zu $\mathscr{S}^n \varphi$ für alle $n \in \Z$.
Aufgrund von 
\begin{align*}
\langle \mathscr{S}^n \varphi, 1 - \varphi \rangle 
= \int \limits_\mathbb{T} \mathscr{S}^n \varphi \overline{(1 - \varphi)} \td{\lambda_N}
= \int \limits_\mathbb{T} (\varphi - | \varphi|^2) z^n \td{\lambda_N}
= 0
\end{align*}
erhalten wir $\varphi = | \varphi|^2$ und finden ein $e \subset \mathbb{T}$ mit $\varphi = \chi_e$.
Da durch $\lbrace z^n \ | \ n \in \Z \rbrace $ ein Basis von $\L^2$ gegeben ist, erhalten wir mit
\begin{align*}
\bigvee_{n \in \Z} \mathscr{S}^n \chi_e :=
\overline{\lin} \bigcup_{n \in \Z} \lbrace \mathscr{S}^n \chi_e \rbrace
= \chi_e \cdot \L^2,
\end{align*}
dass $\chi_e \L^2 \subset E$ gilt.
Wir müssen also noch die Gleichheit zeigen.
Hierfür wählen wir $\alpha \in E \ominus \chi_e \L^2 $.
Damit gilt 
\begin{align*}
\langle 1- \chi_e , \mathscr{S}^n \alpha \rangle = 
\int \limits_\mathbb{T} (1- \chi_e) \overline{\alpha} z^{-n} \td{\lambda_N} = 0
\end{align*}
für $n \in \Z$. Damit folgt $(1- \chi_e) \overline{\alpha} = 0$.
Analog folgt, dass 
\begin{align*}
\forall_{n \in \Z}  \ : \ \langle \alpha , \mathscr{S}^n \chi_e \rangle = 0
\quad \Rightarrow \quad \overline{\alpha} \chi _e= 0
\end{align*}
gilt. Wegen 
\begin{align*}
(1- \chi_e) \overline{\alpha} =
\overline{\alpha} \chi _e= 0
\end{align*}
erhalten wir direkt $\alpha = 0$.
Damit gilt $E = \chi_e \L^2$.

\item
In diesem Fall ist $\mathscr{S}E$ ein echter Teilraum von $E$.
Wir wählen ein $\Theta \in E \ominus \mathscr{S} E$ mit $\| \Theta \|_2 = 1$.
Der Operator $\mathscr{S}$ ist unitär, weswegen wir
\begin{align*}
\langle \mathscr{S}^n \Theta , \Theta \rangle
= 
\langle \Theta, \mathscr{S}^{-n} \Theta \rangle
=
\int \limits_\mathbb{T} z^n | \Theta |^2 \td{\lambda_N} =0 
\end{align*}
für $n \in \N$ erhalten.
Damit ist $|\Theta| $ konstant.
Da $\lambda_N(\mathbb{T}) = 1$ gilt, folgt mit 
\begin{align*}
\langle \Theta, \Theta \rangle 
=
\int \limits_\mathbb{T} | \Theta|^2 \td{\lambda_N} = 1 
\end{align*}
auch $| \Theta  | = 1 $ fast überall.
Wir betrachten $\Theta$ nun als Multiplikationsoperator, d.h. $f \mapsto \Theta f$. 
Damit ist $\Theta$ eine Isometrie auf $\L^2$.
Mit der Einbettung
\begin{align*}
\iota \ : \H^2 \hookrightarrow \L^2, \quad 
z^n \mapsto z^n
\end{align*}
und der Isometrie erhalten wir durch
\begin{align*}
\bigvee_{n \in \N_0} \mathscr{S}^n \Theta
=
\Theta \bigvee_{n \in \N_0} \mathscr{S}^n 1 
=
\Theta \bigvee_{n \in \N_0} \iota(z^n) 1
= \Theta H^2, 
\end{align*}
dass $\Theta \H^2 \subset E$ gilt.
Übrig zu zeigen ist wieder die Gleichheit.
Wir wählen $\alpha \in E \ominus \Theta \H^2$.
Dann gilt 
\begin{align*}
\langle \alpha , \mathscr{S}^n \Theta \rangle
= 
\langle \mathscr{S}^{-n} \alpha , \Theta \rangle
= 0
\end{align*} 
für $n \in \N_0$ und wegen $ \mathscr{S}^n \alpha \in \mathscr{S} E$ gilt auch
\begin{align*}
\langle \Theta, \mathscr{S}^n \alpha \rangle 
=
\langle \mathscr{S}^{-n} \Theta, \alpha \rangle
= 0
\end{align*}
für $n \in \N$.
Damit erhalten wir
\begin{align*}
\int \limits_\mathbb{T} \Theta \overline{\alpha} z^n \td{\lambda_N} = 0
\end{align*}
für alle $n \in \Z$.
Damit folgt direkt $\alpha = 0$.
Somit gilt $E = \Theta H^2$.
\end{enumerate}
\end{proof} 

\begin{lemma}\label{skript:1.2}
Sei $E$ ein Unterraum von $\L^2$.
Es gelten:
\begin{enumerate}[label=\textbf{(\arabic*)}]
\item
Jeder Unterraum der Form $\chi_e \L^2$ ist $\mathscr{S}$ und $\mathscr{S}^{-1}$-invariant.
Die Darstellung $E = \chi_e \L^2$ ist eindeutig.

\item
Jeder Unterraum der Form $\Theta H^2$, $| \Theta | = 1$ f.ü. ist $\mathscr{S}$-invariant, jedoch nicht $\mathscr{S}^{-1}$-invariant.
Die Darstellung $E = \Theta H^2$ ist bis auf eine multiplikative Konstante eindeutig. 

\end{enumerate}
\end{lemma}

\begin{proof}
\begin{enumerate}[label=\textbf{(\arabic*)}]
\item
Folgt direkt aus dem letzten Beweis.
\item
Die Invarianz bzgl. $\mathscr{S}$ und Nichtinvarianz bzgl. $\mathscr{S}^{-1}$ folgen auch hier direkt aus dem letzten Beweis.
Wir betrachten nun $\Theta \H^2 = \Theta^\prime \H^2$.
Dann folgt $\Theta \overline{\Theta^\prime} \H^2 = H^2$
bzw $\H^2 = \Theta^\prime \overline{\Theta} \H^2$.
Nun gilt $\Theta \overline{\Theta^\prime} \in \H^2$ und $\Theta \overline{\Theta^\prime} \in \overline{\H^2}$.
Seien $f , \overline{f} \in \H^2$.
Dann erhalten wir 
\begin{align*}
\hat{f}(n) = 0
\end{align*}
und
\begin{align*}
\overline{\hat{f}}(n) = \overline{\hat{f}(-n)} = 0
\end{align*}
für $n < 0$. Somit ist $f$ konstant und wir erhalten die Eindeutigkeit bis auf eine multiplikative Konstante.
%Mit
%\begin{align*}
%\bigvee_{n\in \N_0} \mathscr{S}^n \Theta 
%= 
%\bigvee_{n \in \N_0} \mathscr{S}^n \Theta^\prime
%\quad \Rightarrow \quad
%\Theta \mathscr{S}^n 1 = C \cdot \Theta^\prime \mathscr{S}^n 1
%\end{align*}
%folgt die restliche Aussage sofort. 
\end{enumerate}
\end{proof}

\begin{genericthm}{Folgerung}\label{skript:1.3}
Sei $E \subset \H^2$ ein nicht trivialer Unterraum und $\mathrm{S}E \subset E$.
Dann gilt $E = \Theta \H^2$ mit $| \Theta | = 1$ fast überall auf $\mathbb{T}$.
\end{genericthm}

\begin{proof}
Mit der Einbettung 
\begin{align*}
\iota \ : \H^2 \hookrightarrow \L^2, \quad 
z^n \mapsto z^n
\end{align*}
gilt $\iota(E) \subset \iota(\H^2) \subset \iota(\L^2)$ und $\mathscr{S} E  \subset E$.
Aus Beurling-Helson erhalten wir sofort
\begin{align*}
E = \iota(E) = \Theta \H^2
\end{align*}
mit $|\theta | = 1$ f.ü. auf $\mathbb{T}$.
Wegen $E \subset \H^2$ gilt auch $\Theta \in \H^2$.
\end{proof}


\begin{genericthm}{Folgerung}\label{skript:1.4}
Sei $f \in \H^2$ und $\lambda_N( \lbrace \xi \in \mathbb{T}\ | \ f(\xi) = 0 \rbrace) > 0$.
Dann gilt $f = 0$.
\end{genericthm}

\begin{proof}
Zunächst ist klar, dass
\begin{align*}
A:=\lbrace \xi \in \mathbb{T}\ | \ f(\xi) = 0 \rbrace
\end{align*}
bis auf eine Nullmenge eindeutig ist. Wir betrachten nun $f$ als Element des $\L^2$.
Dann ist 
\begin{align*}
E_f = \bigvee_{n \in \N_0} \mathscr{S}^n f
\end{align*}
eine $\mathscr{S}$-invarianter Unterraum.
Die in $E_f$ enthaltenen Funktionen verschwinden auf $A$.
Aus $f \neq 0$ folgt zunächst $E_f \neq {0}$.
Mit \ref{skript:1.3} erhalten wir 
$E_f = \Theta \H^2$ mit $\Theta \in E_f$ und $|\Theta | = 1$ fast überall.
\end{proof}
\newpage
\section{Faktorisierung}

\begin{df}\label{skript:2.1}
	Sei $ f \in \H^2 $.
	Wir bezeichnen $ f $ als \textit{innere Funktion}, falls $ | f | = 1 $ fast überall auf $ \mathbb{T} $ gilt.
	Wir nennen $ f $ eine \textit{äußere Funktion}, falls
	\begin{align*}
	E_f := \bigvee_{n \in \N_0} \mathscr{S}^n f = \H^2
	\end{align*}
	gilt.
\end{df}

\begin{genericthm}{Faktorisierungstheorem}\label{skript:2.2}
	Sei $ f \in \H^2 \setminus \lbrace 0 \rbrace $.
	Dann lässt sich $ f $ bis auf eine multiplikative Konstante eindeutig als Produkt einer inneren Funktion $ f^i $ und einer äußeren Funktion $ f^a $ beschreiben.
	D.h. es gilt $ f = f^i f^a $ und $ E_f = f^i \H^2 $.
\end{genericthm}

\begin{proof}
	Wegen $ f \neq 0 $ erhalten wir  mit \ref{skript:1.3}
	\begin{align*}
	E_f = \Theta \H^2	
	\end{align*}
	mit $ |\Theta| = 1  $ f. ü. auf $ \mathbb{T} $.
	Damit ist $ \Theta  $ eine innere Funktion. Wir setzen $ f^i := \Theta $ und $ f^a := f \overline{\Theta} $, womit gilt $ f = f^i f^a $. 
	Wir müssen noch nachweisen, dass $ f^a $ eine äußere Funktion ist.
	Da $ f^i $ als Multiplikationsoperator eine Isometrie auf $ \L^2 $ bzw. $ \H^2 $ ist, erhalten wir 
	\begin{align*}
	E_f
	= \bigvee_{n \in \N_0} z^n f^i f^a 
	= f^i \bigvee_{n \in \N_0} z^n f^a 
	= f^i E_{f^a}.
	\end{align*}
	Damit folgt $ f^i H^2 = f^i E_{f^e} $ und damit auch $ H^2 = E_{f^e} $. Also ist $ f^e $ eine äußere Funktion.
	Die Eindeutigkeit erhalten wir aus Beurling-Helson und \ref{skript:1.2}.
\end{proof}

\begin{genericthm}{Folgerung}\label{skript:2.3}
	Sei $ f \in \H^2$.
	Dann gilt:
	\begin{align*}
	f \ \text{ist äußere Funktion} \
	\Leftrightarrow
	\left(g \in \H^2, \frac{g}{f} \in \L^2 \Rightarrow \frac{g}{f} \in \H^2 \right)
	\end{align*}
\end{genericthm}

\begin{proof}
	Zunächst beginnen wir mit der Hinrichtung.
	Sei $ f $ eine äußere Funktion, $ g \in \H^2$ und $ \frac{g}{f} \in \L^2 $. Wegen $ E_f = \H^2 $ existiert eine Folge $ p_n $ von Polynomen, sodass
	\begin{align*}
	\lim\limits_{n \to \infty} \| p_n f - 1 \|_2 = 0
	\end{align*}
	gilt. Damit erhalten wir 
	\begin{align*}
	\left\| p_n g - \frac{g}{f} \right\|_1 
	= 
	\left\| (p_n f - 1 ) \frac{g}{f} \right\|_1 
	\leq
	\| p_n f - 1 \|_2 \left\| \frac{g}{f} \right\|_2
	\stackrel{n \to \infty}{\rightarrow} 0
	\end{align*}
	mit der Hölderungleichung. Damit gilt
	\begin{align*}
	\left(\frac{g}{f}\right)^{\hat{}}(n) = 0
	\end{align*}
	für $ n <0 $. Somit gilt $ \frac{g}{f} \in \H^2$.\\
	Nun kommen wir zur Rückrichtung.
	Wir setzen $ g = f^a $.
	Dann gilt
	\begin{align*}
	\frac{g}{f} = \frac{1}{f^i} = \overline{f^i} \in \L^2
	\end{align*}
	und mit der Voraussetzung folgt $ \overline{f^i} \in \H^2 $.
	Wegen $ f^i \in \H^2 $ folgt, dass $ f^i $ konstant ist.
	Somit ist $ f $ eine äußere Funktion.
\end{proof}

\begin{genericthm}{Folgerung}\label{skript:2.4}
	\begin{enumerate}[label=\text{\textbf{(\arabic*)}}]
		\item
		Seien $ f$ und $g $ äußere Funktionen mit $ |f| = |g| $ fast überall.
		Dann folgt $ f = \lambda g $ mit $ \lambda \in \mathbb{T} $.
		
		\item Sei $ f \in \L^2 $ und $ \mathscr{S}^{-1} E_f \neq  E_f$.
		Dann existiert eine äußere Funktion $ g $ mit $ |f | = |g|  $ fast überall auf $ \mathbb{T} $.
		
%		\item
%		Falls für $ h \in \L^2 $ mit $ h >0 $ 
%		\begin{align*}
%		\int \limits_{\mathbb{T} } \log h > - \infty
%		\end{align*}
%		gilt, folgt $ |h| = |f| $ fast überall für eine äußere Funktion $ f $.
	\end{enumerate}
\end{genericthm}

\begin{proof}
	\begin{enumerate}[label=\text{\textbf{(\arabic*)}}]
		\item 
		Seien $ f $ und $ g $ äußere Funktionen mit $ |f| = |g| $ fast überall.
		Dann ist mit \ref{skript:2.3}
		\begin{align*}
		\frac{g}{f} \in \H^2 \ \wedge \overline{ \left(\frac{f	}{g} \right)} \in \H^2,
		\end{align*}
		womit folgt, dass $ \frac{g}{f} $ konstant ist. Somit gilt die Aussage.
		\item
		Wir wenden Beurling-Helson an. Damit existiert ein $ \Theta \in \L^2 $ mit $ |\Theta| = 1 $ fast überall, sodass
		\begin{align*}
		E_f = \Theta \H^2
		\ \Leftrightarrow \ 
		\overline{\Theta} E_f = E_{\overline{\Theta} f} = \H^2
		\end{align*}
		gilt. Damit wissen wir, dass ein $ g \in \H^2 $ mit
		\begin{align*}
		\iota(g) = \overline{\Theta} f
		\end{align*}
		existiert. Die Einbettung $ \iota $ ist eine Isometrie, wodurch die Aussage direkt mit \ref{skript:2.2} folgt.
	\end{enumerate}
\end{proof}

\begin{genericthm}{Folgerung}\label{skript:2.5}
	Sei $ f \in \L^2 $. Dann gelten:
	\begin{enumerate}[label=\text{\textbf{(\arabic*)}}]
		\item 
		$ \mathscr{S} E_f \neq E_f \ \Leftrightarrow  \ \log|f|  \in \L^1$
		\item 
		$ E_f = \L^2   \ \Leftrightarrow \ \log|f| \notin \L^1$ mit $ f \neq 0 $ f.ü.
	\end{enumerate}
\end{genericthm} 

\begin{proof}
	\begin{enumerate}[label=\text{\textbf{(\arabic*)}}]
		\item 
		Sei $ \mathscr{S} E_f \neq E_f $. Mit dem selben Argument wie im zweiten Teil der letzten Folgerung erhalten wir
		\begin{align*}
		\H^2 = E_{ \overline{\Theta} f}
		\end{align*}
		mit $ \overline{\Theta} f \in \iota(\H^2) $.
		Dies ist äquivalent zu $ \log| \overline{\Theta} f | = \log |f| \in \L^1 
		$.
		\item 
		Wegen $ E_f = \L^2 $ gilt $ \mathscr{S} E_f = E_f $, womit direkt $ \log|f| \notin \L^1 $ folgt.  
		Wegen 
		\begin{align*}
		E_f \subset \chi_e \L^2, \quad 
		e = \lbrace \xi \in \mathbb{T}\ | \ f(\xi) \neq 0 \rbrace
		\end{align*}
		gilt $ f\neq 0 $ fast überall.
		Anderseits gilt $ \log |f| \notin \L^1 $ mit der ersten Aussage $ \mathscr{S} E_f = E_f $. Beurling-Helson liefert dann $ E_f = \chi_e \L^2 $.
		Mit $ \lambda_N(e) = 1 $ folgt dann $ E_f = \L^2 $.
 	\end{enumerate}
\end{proof}

\begin{df}\label{skript:2.6}
	Eine innere Funktion $ \Theta_1 $ \textit{teilt} die innere Funktion $ \Theta_2 $, falls
	\begin{align*}
	\frac{\Theta_2}{\Theta_1} \in \H^2
		\end{align*}
		gilt. Damit ist $ \frac{\Theta_2}{\Theta_1} $ eine innere Funktion.
		Wir nennen $ \Theta $ das kleinste gemeinsame Vielfache ($ \kgv $) von $ \Theta_1 $ und $ \Theta_2 $, falls $ \Theta $ von $ \Theta_1 $ und $ \Theta_2 $ geteilt wird und jede innere Funktion $ \Theta^\prime $, welche von $ \Theta_1 $ und $ \Theta_2 $ geteilt, auch von $ \Theta $ geteilt wird.
		
\end{df}

\begin{genericthm}{Satz}\label{skript:2.7}
	
	\begin{enumerate}[label=\text{\textbf{(\arabic*)}}]
		\item Es gilt $ \Theta_2 \H^2 \subset \Theta_1 \H^2  $ genau dann, wenn
		$ \Theta_1$ teilt $ \Theta_2 $ gilt.
		
		\item
		Es gilt
		\begin{align*}
		\Theta_1 \H^2 \cap \Theta_2 \H^2 = \Theta \H^2,
		\end{align*}
		wobei $ \theta $ das kleinste gemeinsame Vielfache von $ \Theta_1 $ und $ \Theta_2$ ist.
	\end{enumerate}
	
\end{genericthm}

\begin{proof}
	\begin{enumerate}[label=\text{\textbf{(\arabic*)}}]
		\item
			Wegen $ \Theta_2 \H^2 \subset \Theta_1 \H^2 $ existiert eine innere Funktion $ f $ mit $ \Theta_2 = \Theta_1 f $.\\
			Aufgrund von
			\begin{align*}
			\Theta_2 \H^2 = \Theta_1 f \H^2 \subset \Theta_1 \H^2
			\end{align*}
			erhalten wir $ \Theta_2 \in \Theta_1 \H^2 $. Somit gilt auch 
			$ \frac{\Theta_2}{\Theta_1} \in \H^2 $.
		\item 
			Der Raum $ \Theta_1 \H^2 \cap \Theta_2 \H^2 $ ist $ S $-invariant ist. Des weiteren ist $\Theta= \Theta_1 \Theta_2 \neq 0 $ in diesem Unterraum enthalten.
			Demnach gilt $\Theta_1 \H^2 \cap \Theta_2 \H^2 = \Theta \H^2  $
			Es folgt direkt, dass $ \Theta $ von $ \Theta_1 $ und $ \Theta_2 $ geteilt wird.
			Wenn nun ein $ \Theta^\prime $ von $ \Theta_1$ und $ \Theta_2 $ teilbar ist, erhalten wir mit
			\begin{align*}
			\Theta^\prime \H^2 \subset \Theta_1 \H^2 \ \wedge \ \Theta^\prime \H^2 \subset \Theta_2 \H^2
			\end{align*}
			sofort $ \Theta^\prime \H^2 \subset \Theta \H^2 $.
			Damit ist $ \Theta  $ das kleinste gemeinsame Vielfache von $ \Theta_1 $ und $ \Theta_2 $.
	\end{enumerate}
\end{proof}

 

\newpage
\section{Eigenschaften des adjungierten Operators $ S^\ast $}

\begin{sz}
	\begin{enumerate}[label=\text{\textbf{(\arabic*)}}]
		\item 
		Sei $ E $ ein abgeschlossener Teilraum von $ \H^2 $ mit $ E \neq \H^2 $ und $ S^\ast E \subset E $.
		Dann gilt 
		\begin{align*}
		E = \H^2 \ominus \Theta \H^2
		\end{align*}
		mit einer inneren Funktion $ \Theta $.
		Außerdem ist jeder Unterraum der Form $ \H^2 \ominus \Theta \H^2 $
		$ S^\ast $-  invariant.	
		
		\item
		Wenn $ E $ ein $ S $-und $ S^\ast $-invarianter Unterraum von $ \H^2 $ ist, dann gilt entweder $ E = {0} $ oder $ E = \H^2 $.
	\end{enumerate}
\end{sz}

\begin{proof}
	\begin{enumerate}[label=\text{\textbf{(\arabic*)}}]
		\item 
		Mit $ S^\ast E \subset E $ gilt
		\begin{align*}
		\langle x , S^\ast y \rangle =
		\langle S x ,  y \rangle = 0, \quad x \in E^\bot, y \in E,
		\end{align*}
		wodurch wir sofort
		\begin{align*}
		S^\ast E \subset E \Leftrightarrow S E^\bot \subset E^\bot
		\end{align*}
		erhalten. Mit Beurling-Helson erhalten wir dann $ E^\bot = \Theta \H^2 $, womit die Aussage gilt.
		
		\item
		Angenommen $ \lbrace 0 \rbrace \neq E \neq \H^2 $.
		Wegen $ S E \subset E  $ und $ S^\ast E \subset E $, gelten 
		\begin{align*}
		E = \Theta \H^2 = \H^2 \ominus \Theta^\prime \H^2
		\end{align*}
		für zwei innere Funktionen $ \Theta, \Theta^\prime$.
		Dann gilt 
		\begin{align*}
		\langle \Theta f , \Theta^\prime g \rangle   
		=
		\int \limits_{T} \Theta f \overline{\Theta^\prime g}  \td{\lambda_N} = 0
		\end{align*}
		für $ f, g \in \H^2 $. Für $ \langle \Theta \Theta^\prime, \Theta^\prime \Theta \rangle = 1 $ ist dies ein Widerspruch.		
		\end{enumerate}
\end{proof}

\begin{genericthm}{Folgerung}
	Sei $ \mathcal{F}$  eine Familie innerer Funktionen und 
	$ K_\Theta := \H^2 \ominus \Theta \H^2$ für eine innere Funktion $ \Theta $.
	Dann gelten:
	\begin{enumerate}[label=\text{\textbf{(\arabic*)}}]
		\item $ K_{\Theta_1} \subset K_{\Theta_2}  \ \Leftrightarrow \ \Theta_1 \ \text{teilt} \ \Theta_2$
		
		\item
		$ \bigvee_{\Theta \in \mathcal{F}} K_\Theta  = K_{\Theta^\prime}$ mit $ \Theta^\prime = \kgv(\mathcal{F})  $
		
		\item
		$ \bigcap_{\Theta \in \mathcal{F}} K_\Theta = K_{\Theta^\prime} $
		mit $ \Theta^\prime = \ggT(\mathcal{F}) $.
	\end{enumerate}
	Für den Fall $ \bigvee_{\Theta \in \mathcal{F}} K_\Theta  = \H^2 $ 
	und $ \bigcap_{\Theta \in \mathcal{F}} K_\Theta = \lbrace 0 \rbrace $ sagen wir $ \kgv( \mathcal{ F}) = 0 $ und $ \ggT(\mathcal{F}) = 1 $.
\end{genericthm}

\begin{proof}
	\begin{enumerate}[label=\text{\textbf{(\arabic*)}}]
		\item
		Durch
		\begin{align*}
		K_{\Theta_1} =  (\Theta_1 \H^2)^\bot  \subset (\Theta_2 \H^2)^\bot = K_{\Theta_2}
		\ \Leftrightarrow \ 
		\Theta_2 \H^2 \subset \Theta_1 \H^2
		\ \Leftrightarrow \
		\Theta_1 \ \text{teilt} \ \Theta_2
		\end{align*}
		erhalten wir die Aussage sofort.
		
		\item
		Der Raum 
		\begin{align*}
		\bigvee_{\Theta \in \mathcal{F}} K_\Theta
		\end{align*}
		ist $ S^\ast $-invariant.
		Damit existiert eine innere Funktion $ \Theta^\prime $ mit
		\begin{align*}
		\bigvee_{\Theta \in \mathcal{F}} K_\Theta = K_{\Theta^\prime}
		\end{align*}
		und es gilt
		\begin{align*}
		K_\Theta \subset K_{\Theta^\prime}
		\end{align*}
		für alle $ \Theta \in \mathcal{F} $.
		Damit gilt auch $ \Theta^\prime = \kgv(\mathcal{F}) $.
		
		\item
		Wegen der $ S^\ast $-Invarianz gilt auch hier 
		\begin{align*}
		\bigcap_{\Theta \in \mathcal{F}} K_\Theta = K_{\Theta^\prime}
		\end{align*}
		für eine innere Funktion $ \Theta^\prime $.
		Aus 
		\begin{align*}
		K_{\Theta^\prime} \subset K_{\Theta}
		\end{align*}
		für $ \Theta \in \mathcal{F} $ folgt direkt $ \Theta^\prime = \ggT(\mathcal{F}) $.
	\end{enumerate}
\end{proof}
\end{document}